\chapter{The Input System}
\label{ch:input-system}

CNA's \cnans{Microsoft::Xna::Framework::Input} namespace covers keyboard, mouse, gamepad,
and touch --- and, layered underneath it, a set of CNA-only extensions with no XNA
counterpart at all. The current source audit counts 522 input test cases and a compile-time
strict-API freeze. That is substantial automated evidence, but it does not replace physical
keyboard, mouse, gamepad, touchscreen, IME, and high-DPI verification.

\section{Architecture: event-driven, not poll-driven}

FNA queries SDL when \texttt{GetState()} is called. CNA caches state from the events drained by
\texttt{Game::PollEvents()}. \texttt{SdlInputBridge::ProcessEvent} updates
\texttt{InputManager}, \texttt{TextInputEXT}, and the touch gesture detector; the public
\texttt{GetState()} methods read those snapshots. \texttt{FrameworkDispatcher::Update()} drives
the per-frame touch gesture update.

Repeated reads without another event return the same keyboard, absolute-mouse, gamepad, and touch
state. Relative mouse mode is the exception: the first \texttt{Mouse::GetState()} drains its
motion accumulator, and a second read returns zero until more motion arrives. The event pump and
state reads are single-threaded and unsynchronized; both belong on the game-loop/video thread.

Two ordinary keys have an additional framework side effect: \texttt{F9} and \texttt{F10}
simulate graphics-context loss and restoration in \texttt{Game::PollEvents()}. The event is sent
to \texttt{SdlInputBridge::ProcessEvent} first, so both keys still reach
\texttt{Keyboard::GetState()}; application code may observe them but cannot suppress the debug
action. Avoid binding gameplay to them unless that game-loop behavior is changed.

Three tags recur in the headers: \textbf{STRICT} (the frozen FNA/XNA-shaped surface),
\textbf{EXT} (an FNA-compatible extension, using the \texttt{...EXT} suffix convention from
this book's game-loop chapter), and \textbf{CNAEXT} (a CNA-only addition with no FNA counterpart
at all). \texttt{PublicApiInputSignatureFreezeTests.cpp} compile-checks the strict member set.

\section{Keyboard}

\cnaclass{Keyboard} is static-only, offering \texttt{GetState()} and a
per-\cnaclass{PlayerIndex} overload. \cnaclass{KeyboardState} stores keys in a
\texttt{std::unordered\_set<Keys>} instead of FNA's bitfield while preserving the observable
queries, sorted \texttt{GetPressedKeys()} result, and FNA-shaped hash formula. \cnaclass{Keys} has
exactly
160 values, independently confirmed byte-for-byte identical to FNA's numeric values,
including several distinctly odd hexadecimal outliers (\texttt{Pause = 0x13},
\texttt{Kana = 0x15}, \texttt{ChatPadGreen = 0xCA}) preserved exactly rather than renumbered
for tidiness. The SDL bridge has 125 scancode-to-key cases but 126 cases in the reverse
direction, and many enum values have no SDL scancode source at all. IME keys, ChatPad keys, and
several browser/media keys therefore cannot round-trip through the \texttt{EXT} scancode
helpers. Typed text must go through \texttt{TextInputEXT}, not through assumptions that every
Unicode character has a \cnaclass{Keys} value.

\subsection{KeyboardState and keyboard extensions}

\cnaclass{KeyboardState}'s public surface is compact:

\begin{description}
\item[\texttt{KeyboardState(std::initializer\_list<Keys> keys)}] the ordinary way application
  and test code constructs one, for example
  \texttt{KeyboardState\{Keys::W, Keys::LeftShift\}}.
\item[\texttt{getItem(Keys key)} / \texttt{operator[](Keys key)}] both return a
  \cnaclass{KeyState} (\texttt{Down} or \texttt{Up}) for one key.
\item[\texttt{IsKeyDown(key)} / \texttt{IsKeyUp(key)}] boolean key queries.
\item[\texttt{GetPressedKeys()}] returns every currently-down key as an ascending-sorted
  \texttt{std::vector<Keys>}; the sort is part of the contract because the backing hash set has
  no stable iteration order.
\end{description}

\cnaclass{Keyboard}'s \texttt{EXT} methods bridge XNA's US-QWERTY-shaped \cnaclass{Keys} enum
and SDL3's distinction between a physical scancode and a layout-dependent keycode.
\texttt{GetKeyFromScancodeEXT()} maps a US-named physical position into the current layout,
which supports position-based movement controls on non-US keyboards. The public API has no
general keycode-to-physical-position inverse; the inverse helpers apply to textual names.
\texttt{GetModStateEXT()} returns SDL's modifier mask. \texttt{GetScancodeNameEXT} and
\texttt{GetKeyNameEXT}, plus their name-to-key forms, provide current-layout labels for rebinding
interfaces.

\begin{lstlisting}
// Physical-position movement keys for a non-QWERTY layout.
// WASD's physical positions, expressed as scancodes rather than as the
// QWERTY-specific Keys::W/A/S/D values directly.
const Keys forwardKey = Keyboard::GetKeyFromScancodeEXT(Keys::W);
const Keys leftKey    = Keyboard::GetKeyFromScancodeEXT(Keys::A);
const Keys backKey     = Keyboard::GetKeyFromScancodeEXT(Keys::S);
const Keys rightKey   = Keyboard::GetKeyFromScancodeEXT(Keys::D);

// On an AZERTY keyboard, forwardKey resolves to the physical key labelled Z.
// GetKeyFromScancodeEXT(Keys::W) treats Keys::W as identifying the
// physical scancode a US layout would call W, then maps that scancode
// back through the *current* layout to find what key is physically there.
const KeyboardState keyboard = Keyboard::GetState();
if (keyboard.IsKeyDown(forwardKey))
{
    // move forward
}
\end{lstlisting}

\section{Mouse}

\cnaclass{Mouse} is static-only: \texttt{GetState()} and \texttt{SetPosition(x, y)}, the
latter converting logical coordinates to physical window coordinates before calling
SDL's warp function. The shared routing first asks an associated SDL renderer, then a
registered graphics renderer, then falls back to pass-through. Those tiers make the result
renderer- and presentation-dependent; the audit below and
Chapter~\ref{ch:backend-architecture}'s matrix state the exact boundary.
\cnaclass{MouseState} carries the expected X/Y, five buttons, and a cumulative,
120-units-per-notch \texttt{ScrollWheelValue}; a \texttt{CNAEXT} horizontal-scroll addition is
deliberately excluded from \texttt{Equals} / \texttt{GetHashCode} / \texttt{ToString} specifically
so those stay byte-identical to FNA's own output despite the extra field existing. A wholly
\texttt{CNAEXT} class, \cnaclass{MouseCursor}, wraps \texttt{SDL\_Cursor*} with eleven lazily
created, process-lifetime stock system cursors. MonoGame's
\texttt{PlatformDispose()} frees a shared stock cursor handle unconditionally, which can
invalidate other users; CNA makes \texttt{Dispose()} a no-op for shared stock instances.

\subsection{MouseState and mouse extensions}

\cnaclass{MouseState} is a plain data carrier --- \texttt{X} / \texttt{Y}, five
\cnaclass{ButtonState} properties (\texttt{LeftButton} / \texttt{RightButton} /
\texttt{MiddleButton} / \texttt{XButton1} / \texttt{XButton2}), the cumulative
\texttt{ScrollWheelValue} already named above, and (\texttt{CNAEXT}) a parallel
\texttt{HorizontalScrollWheelValueEXT} for horizontal scroll, deliberately excluded from
\texttt{Equals} / \texttt{GetHashCode} / \texttt{ToString} as already noted. \cnaclass{Mouse}
itself, beyond \texttt{GetState()} and \texttt{SetPosition}, exposes relative
(pointer-lock-style) mouse mode and window/capture-related \texttt{EXT}
methods:

\begin{description}
\item[\texttt{IsRelativeMouseModeEXT}] (get/set)
  toggles SDL's own relative mouse mode, where motion is reported as unbounded deltas rather
  than an absolute, screen-clamped position --- the standard technique for an FPS-style camera
  look. Both accessors are safe to call with no window at all (no published window handle and
  no focused window): the getter simply reports \texttt{false} rather than querying SDL with a
  null window (whose behavior SDL leaves undefined), and the setter is a no-op. While relative
  mode is active, each
  \texttt{SDL\_EVENT\_MOUSE\_MOTION} adds its relative delta to a separate accumulator;
  motion outside that mode is not accumulated, and toggling the mode flushes stale deltas.
  \texttt{MouseState}'s public \texttt{X} / \texttt{Y} read and immediately drain that
  accumulator in relative mode; otherwise they report the absolute position.
\item[\texttt{SetPosition(x, y)}] is a documented no-op in relative mode because an absolute
  warp has no meaning while the mouse reports deltas.
\item[\texttt{SetCaptureEXT(enabled)}] engages OS-level mouse capture (input keeps arriving
  even if the cursor leaves the window bounds), distinct from relative mode and combinable
  with it.
\item[\texttt{GetGlobalPositionEXT} / \texttt{WarpGlobalEXT}] read or set the mouse position in
  \emph{desktop} rather than window-relative coordinates.
\item[\texttt{ClickedEXT}] a \texttt{CNAEXT System::MulticastAction<int>} event (Chapter~\ref{ch:sharp-runtime-namespaces}
  covers this \texttt{sharp-runtime} type) firing once per click with the button
  index, for code that wants a discrete click event rather than polling button state every
  frame.
\end{description}

\begin{lstlisting}
// FPS-style camera input: relative mode reports motion as deltas.
game.setIsMouseVisibleProperty(false);
Mouse::setIsRelativeMouseModeEXTProperty(true);

// Each frame: read accumulated relative motion, not an absolute position.
const MouseState mouse = Mouse::GetState();
const float yaw   = mouse.X * mouseSensitivity;
const float pitch = mouse.Y * mouseSensitivity;
cameraRotation = Quaternion::CreateFromYawPitchRoll(yaw, pitch, 0.0f) * cameraRotation;
\end{lstlisting}

\subsection{Logical coordinates are a renderer contract}
\label{sec:input-logical-coordinate-routing}

The write direction is only half the story. Incoming
\texttt{SDL\_EVENT\_MOUSE\_MOTION} and mouse-button positions pass through
\texttt{SdlInputBridge::to\_logical\_position()} before updating
\texttt{InputManager}; touch starts from normalized window position and reaches the same
helper. Thus a bad or missing transform affects absolute mouse reads, click coordinates, and
touch together, while relative mouse motion deliberately remains an unscaled delta.

The SDL\_RENDERER route has the strongest existing proof. A test creates a
$100\times100$ logical presentation inside a $200\times200$ window, injects a physical motion
event at $(100,100)$, and observes logical $(50,50)$ through
\texttt{Mouse::GetState()}. Separate conversion tests cover the inverse and a horizontal
letterbox offset. The latter tests cannot inspect the OS cursor under Wayland, however:
\texttt{SetPosition()} writes the requested logical point into \texttt{InputManager}
\emph{before} attempting conversion and warp, so immediately reading that point back proves
the cache assignment even if the physical cursor landed elsewhere. The former ASCII renderer's integration
test uses this same weak immediate set/get shape.

Other renderers in the inherited audit cohort are not interchangeable:

\begin{itemize}
\item EasyGL and Canvas register an offset-free height scale. It matches their default
  fixed-height, dynamic-width mode, but not the other presentation modes they store.
\item WebGPU and SDL GPU implement all five viewport modes, but their transform uses physical
  pixel dimensions for SDL event and warp values expressed in window coordinates. A
  high-pixel-density display therefore introduces a density-factor error. Reads in a
  letterbox bar also return \texttt{false}, causing raw window coordinates to replace the
  outside logical point the method had already computed.
\item Vulkan has no override or registry entry. Its SpriteBatch projection stretches the
  virtual canvas across the physical swapchain, so absolute input becomes misaligned after a
  non-1:1 resize. Keeping its window-coordinate size equal to the requested virtual size is
  the only portable caller-side guard today.
\item BGFX and D3D9/11/12 also pass through, but those renderers in the inherited audit cohort ignore logical
  presentation and draw in physical coordinates; input matches that limitation rather than
  the requested mode.
\item FREEDIRECT's correct direct helper methods are shadowed by the SDL renderer created inside
  \repolink{free-direct}. Public mapping is pass-through until the first
  \texttt{Present()} installs its hardcoded letterbox, then follows SDL's mapping. Its
  four-property test invokes the helpers directly and does not exercise this public route.
\end{itemize}

The complete inherited fourteen-product matrix and its test evidence are in
Chapter~\ref{ch:backend-architecture},
\S\ref{sec:backend-input-coordinate-matrix}. FNA has no corresponding renderer split:
\texttt{Mouse.GetState()} multiplies window X/Y by the independent backbuffer/window ratios,
and \texttt{SetPosition()} applies their inverses. CNA needs more machinery because its
presentation modes include offsets and dynamic logical dimensions, but callers should not
mistake that richer vocabulary for uniform implementation.

\section{GamePad}

The full XNA surface --- \texttt{GetCapabilities}, \texttt{GetState} with configurable
\cnaclass{GamePadDeadZone} handling, \texttt{SetVibration} --- is wired to real
\texttt{SDL\_Gamepad} hardware, not stubbed. CNA exposes at most four player slots; an
out-of-range player-index override is clamped rather than creating a fifth slot. Dead-zone
constants match FNA exactly
(\texttt{LeftDeadZone = 7849/32768}, \texttt{RightDeadZone = 8689/32768},
\texttt{TriggerThreshold = 30/255}). Beyond the strict XNA surface, an extensive
\texttt{EXT} layer exposes real hardware capability FNA itself already extends beyond stock
XNA: GUID, light-bar and trigger-rumble control, gyroscope/accelerometer reads, and several
CNA-only additions with no FNA counterpart at all (player-index assignment, power info,
button-label queries, device name/path/serial/firmware, Steam handle, connection state,
touchpad finger tracking).

\subsection{Bugs found and fixed}

\begin{historynote}
The gamepad route originally omitted \texttt{SDL\_INIT\_GAMEPAD}, so it received no events.
Another implementation queried rumble support by calling \texttt{SDL\_RumbleGamepad(0,0,0)},
which stopped active vibration; it now uses a non-mutating property query. Stick normalization
used 32768 for negative samples instead of FNA's 32767, making $-16384$ yield $-0.5$ instead of
approximately $-0.50001$. Finally, the initialized gamepad subsystem was not shut down on device
recreation. An idempotent \texttt{ShutdownGamepadSubsystem()} now runs during game disposal.
\end{historynote}

\subsection{GamePadState and dead-zone modes}

\cnaclass{GamePadState} composes four sub-structures, each a plain value carrier read directly
from its header: \cnaclass{GamePadThumbSticks} (\texttt{Left} / \texttt{Right}, each a
\cnaclass{Vector2}), \cnaclass{GamePadTriggers} (\texttt{Left} / \texttt{Right}, each a scalar
\texttt{float}), \cnaclass{GamePadButtons} (eleven \cnaclass{ButtonState} properties --- every
face/shoulder/stick-click/back/start/big button --- plus a \texttt{ButtonStateFromFlag(Buttons)}
convenience for looking one up by the \cnaclass{Buttons} flag enum instead of by name), and
\cnaclass{GamePadDPad} (\texttt{Up} / \texttt{Down} / \texttt{Left} / \texttt{Right}, each a
\cnaclass{ButtonState}). \cnaclass{GamePadState} itself adds \texttt{IsConnected},
\texttt{PacketNumber} (incremented when the hardware state changes), and
\texttt{IsButtonDown} / \texttt{IsButtonUp} convenience
wrappers over the same \cnaclass{Buttons} flag enum.

The three \cnaclass{GamePadDeadZone} modes differ in shape as well as threshold.
\texttt{None} filters nothing. The default \texttt{IndependentAxes} excludes and clamps each
axis separately, which can distort direction near the boundary. \texttt{Circular} excludes by
vector magnitude and clamps to the unit circle. \texttt{None} and \texttt{IndependentAxes}
retain a square clamp, so changing modes can also change maximum diagonal magnitude.

\begin{lstlisting}
// Dead-zone-aware movement with a circular clamp and rumble feedback.
const GamePadState pad = GamePad::GetState(PlayerIndex::One, GamePadDeadZone::Circular);
if (pad.getIsConnectedProperty())
{
    const Vector2 moveInput = pad.getThumbSticksProperty().getLeftProperty();
    playerPosition += moveInput * moveSpeed * deltaSeconds;

    if (pad.IsButtonDown(Buttons::A))
    {
        // Independent left/right motor strengths in [0, 1].
        GamePad::SetVibration(PlayerIndex::One, 0.0f, 0.6f);
    }
}
\end{lstlisting}

Beyond the strict surface, \texttt{GetGyroEXT} / \texttt{GetAccelerometerEXT} expose a
controller's built-in motion sensors, distinct from host-device sensors. They require hardware
that reports them, and
\texttt{GamePadCapabilities::getHasGyroEXTProperty()} / \texttt{getHasAccelerometerEXTProperty()}
exist specifically so calling code can check before relying on either. \texttt{SetLightBarEXT}
and \texttt{SetTriggerVibrationEXT} (independent trigger motors rather than the two
whole-controller motors) are modern-controller extensions, gated by their
\texttt{GamePadCapabilities} flags.

\section{Touch and gestures}

\cnaclass{TouchPanel} is static-only (\texttt{MAX\_TOUCHES = 8}), providing
\texttt{GetCapabilities()}, \texttt{GetState()}, and \texttt{ReadGesture()}.
\cnaclass{TouchCollection} is immutable only by contract, not by construction --- its
\texttt{IsReadOnly} flag is advisory, matching an identical inconsistency already present in
FNA's own implementation. \cnaclass{GestureType} is an eleven-value bit-flag enum (Tap,
DoubleTap, Hold, Drag, Flick, Pinch, and their variants).
One value is \texttt{None}; all ten actual gesture types are implemented by
\texttt{GestureDetector}, including both completion events. There is no declared gesture that
silently lacks a detector.

\begin{historynote}
The touch route once ignored \texttt{SDL\_EVENT\_FINGER\_CANCELED}, leaving canceled touches
active. \texttt{GetCapabilities()} used only an observed-touch flag, so an untouched enumerated
device appeared disconnected. Touch-state advancement also occurred inside \texttt{GetState()},
making repeated reads mutate a frame. The current implementation handles cancellation, enumerates
SDL devices with the observed flag as a fallback, and advances state once per frame through
\texttt{AdvanceTouchFrame()}.
\end{historynote}

\subsection{TouchCollection, TouchLocation, and GestureSample}

\cnaclass{TouchCollection} is a per-frame snapshot of tracked touches. It provides
\texttt{Count} / \texttt{IsConnected} (the sticky ``a touch has been observed'' flag described
below), the by-index \texttt{operator[]}, and the same general
ordered-container surface Chapter~\ref{ch:math-core-types}'s own \cnaclass{CurveKeyCollection} coverage already
described --- \texttt{Contains} / \texttt{IndexOf}, a \texttt{CNAEXT} \texttt{empty()}, and
\texttt{CNAEXT} iterator support for range-based \texttt{for}. Each \cnaclass{TouchLocation} inside it carries
an \texttt{Id} stable across frames for one finger, a \texttt{State} (\cnaclass{TouchLocationState}:
\texttt{Pressed}, \texttt{Moved}, \texttt{Released}, or \texttt{Invalid} for a stale/expired
entry), a \texttt{Position}, and (\texttt{CNAEXT}) a \texttt{PressureEXT} reading on hardware
that reports one. \cnaclass{GestureSample} (returned by \texttt{TouchPanel::ReadGesture()})
carries a \texttt{GestureType}, a \texttt{Timestamp}, and up to two positions/deltas
(\texttt{Position} / \texttt{Position2}, \texttt{Delta} / \texttt{Delta2}); the second pair is used
only by two-finger gestures like \texttt{Pinch}, left at their default for single-finger
gestures.

The two connected flags do not use the same evidence. \texttt{TouchPanel::GetCapabilities()}
enumerates SDL touch devices on every call, then falls back to the sticky observed-touch flag or a
currently live touch. It can therefore report a device before the first interaction.
\texttt{TouchCollection::IsConnected}, in contrast, reads only that sticky flag; it can remain
false for an enumerated but untouched device until the first finger-down event. Use capabilities
for device discovery and the collection's count for whether a finger is currently tracked.

\begin{lstlisting}
// Per-finger drag tracking uses the stable Id, not collection position.
const TouchCollection touches = TouchPanel::GetState();
for (const TouchLocation& touch : touches)
{
    switch (touch.getStateProperty())
    {
    case TouchLocationState::Pressed:
        activeDrags[touch.getIdProperty()] = touch.getPositionProperty();
        break;
    case TouchLocationState::Moved:
        if (auto it = activeDrags.find(touch.getIdProperty()); it != activeDrags.end())
        {
            const Vector2 delta = touch.getPositionProperty() - it->second;
            PanCamera(delta);
            it->second = touch.getPositionProperty();
        }
        break;
    case TouchLocationState::Released:
    case TouchLocationState::Invalid:
        activeDrags.erase(touch.getIdProperty());
        break;
    }
}

// Gestures are read separately from raw touches, and only the types
// enabled via EnabledGestures are ever produced by ReadGesture().
TouchPanel::setEnabledGesturesProperty(GestureType::Pinch | GestureType::Flick);
while (TouchPanel::getIsGestureAvailableProperty())
{
    const GestureSample gesture = TouchPanel::ReadGesture();
    if (gesture.getGestureTypeProperty() == GestureType::Pinch)
    {
        // Position/Position2 are the two fingers' current positions;
        // Delta/Delta2 (unused by single-finger gestures) hold their
        // per-frame motion -- combine both to derive a zoom factor.
        ZoomCamera(gesture.getPositionProperty(), gesture.getPosition2Property());
    }
}
\end{lstlisting}

\texttt{EnabledGestures} is an active filter: \texttt{ReadGesture()} returns only enabled
types. \texttt{GestureType} uses powers of two from \texttt{Tap = 1} through
\texttt{PinchComplete = 512}, so callers combine selections with \texttt{|}.

\cnaclass{GestureSample::Timestamp} deliberately differs from an FNA defect. FNA computes it as
\texttt{TimeSpan.FromTicks(Environment.TickCount)} --- but \texttt{TickCount} is a millisecond
value, while \texttt{FromTicks} expects 100-nanosecond units. CNA converts milliseconds to ticks.
Neither engine defines an absolute epoch for the field; the regression establishes non-negative,
strictly increasing values across two gestures under an advanced test clock.

\section{TextInputEXT and CNA-only CNAEXT extensions}

\texttt{TextInputEXT} is not itself a CNA invention --- FNA ships this same extension beyond
the strict XNA~4.0 surface, and CNA ports it directly: \texttt{TextInput} / \texttt{TextEditing}
events, \texttt{StartTextInput} / \texttt{StopTextInput}, on-screen-keyboard visibility, and
input-rectangle placement. The implementation has three event channels: committed
\texttt{TextInput}, in-progress \texttt{TextEditing}, and the CNA-only
\texttt{TextEditingCandidatesEXT} candidate-list event. CNA also adds an explicit
on-screen-keyboard/IME type hint.

The \texttt{CNA::Input} namespace also hosts these \texttt{CNAEXT} subsystems:
\cnaclass{Clipboard} (system clipboard text); \cnaclass{Joysticks} (the raw, unmapped SDL
joystick API --- axes, buttons, hats, trackballs --- deliberately independent of
\cnaclass{GamePad}'s own mapped view of the same physical device); \cnaclass{Sensors} (host
device accelerometer/gyroscope, distinct from a gamepad's own motion sensors covered above);
\cnaclass{Power} (battery state); and \cnaclass{Haptics} (SDL3 force feedback). Haptic actuation
requires capable hardware
\emph{and} an initialized SDL haptic subsystem; most gamepads support only simple
rumble and this layer does not initialize \texttt{SDL\_INIT\_HAPTIC} itself).
This entire \texttt{CNA::Input} extension layer is compiled unconditionally, even when
\texttt{CNA\_DEVICES} and \texttt{CNA\_CNAEXT} are off. \cnaclass{Clipboard} and
\cnaclass{Power} overlap the optional \texttt{CNA::Devices} surface; because
\texttt{CNA\_DEVICES} defaults off, the input versions are the only clipboard and host-power
APIs present in a default build.

\subsection{TextInputEXT contract}

\texttt{TextInput} is a \texttt{System::MulticastAction<charcs>} that fires once per composed
UTF-16 unit; a non-BMP code point arrives as a surrogate pair. \texttt{TextEditing} carries the
uncommitted IME candidate, cursor start, and selection length. \texttt{StartTextInput()} and
\texttt{StopTextInput()} enable or disable this event period, and
\texttt{IsTextInputActive()} reports it. \texttt{SetInputRectangle()} positions a composition
popup in logical window coordinates, while \texttt{IsScreenKeyboardShown()} reports a platform
keyboard. The CNA-only \texttt{StartTextInputWithTypeEXT()} adds a plain-text, numeric, email,
or similar mobile keyboard hint.

\begin{lstlisting}
// An in-game chat box consumes composed text rather than raw key states.
System::MulticastAction<charcs>::Token textInputToken =
    System::MulticastAction<charcs>::InvalidToken;

void ChatBox::Open()
{
    TextInputEXT::SetInputRectangle(chatBoxScreenBounds);
    TextInputEXT::StartTextInput();
    if (textInputToken == System::MulticastAction<charcs>::InvalidToken)
    {
        textInputToken = TextInputEXT::TextInput.Add(
            [this](charcs c) { composedText += c; });
    }
}

void ChatBox::Close()
{
    TextInputEXT::StopTextInput();
    TextInputEXT::TextInput.Remove(textInputToken);
    textInputToken = System::MulticastAction<charcs>::InvalidToken;
    composedText.clear();
}
\end{lstlisting}

\subsection{CNAEXT-only device subsystem examples}

These five subsystems are small static \texttt{CNAEXT} classes with no XNA counterpart.
\cnaclass{Clipboard} wraps SDL with \texttt{GetTextEXT()}, \texttt{SetTextEXT(text)}, and
\texttt{HasTextEXT()}. \cnaclass{Joysticks}
exposes the \emph{raw}, unmapped view of connected joystick hardware
(\texttt{GetJoysticksEXT()} enumerates every connected device;
\texttt{GetCapabilitiesEXT(id)} / \texttt{GetStateEXT(id)} read one device's shape and current
state directly by axis/button/hat index) --- deliberately independent of \cnaclass{GamePad}'s
own mapped, Xbox-360-shaped view of the same physical hardware, for the rarer case where an
application needs a device's true native layout rather than GamePad's normalized
abstraction. \cnaclass{Sensors} (\texttt{GetSensorsEXT()},
\texttt{GetAccelerometerEXT} / \texttt{GetGyroscopeEXT}) reads the \emph{host device}'s own
motion sensors --- a laptop's built-in accelerometer, a Steam Deck's gyro --- distinct from a
gamepad controller's own motion sensors already covered under \texttt{GamePad}'s
\texttt{EXT} surface above. The current implementation never initializes
\texttt{SDL\_INIT\_SENSOR} itself, however. Unless another subsystem has already initialized
SDL sensors, enumeration is empty and both readers return \texttt{false}; fake-backend tests do
not expose this latent production no-op. \cnaclass{Power} is a single method,
\texttt{GetInfoEXT(secondsLeft\&, percent\&)}, returning a \cnaclass{PowerStateEXT} plus two
output parameters --- the host device's own battery state, not a gamepad's (that is
\texttt{GamePad::GetPowerInfoEXT} instead).

\begin{lstlisting}
// Host and controller batteries use separate query methods.
int hostSecondsLeft = 0, hostPercent = 0;
if (Power::GetInfoEXT(hostSecondsLeft, hostPercent) == PowerStateEXT::OnBattery
    && hostPercent < 15)
{
    ShowLowBatteryWarning("this device");
}

int gamepadPercent = 0;
if (GamePad::GetPowerInfoEXT(PlayerIndex::One, gamepadPercent) == PowerStateEXT::OnBattery
    && gamepadPercent < 15)
{
    ShowLowBatteryWarning("controller 1");
}
\end{lstlisting}

\cnaclass{Haptics} maintains device and effect lifetimes. \texttt{GetHapticsEXT()} enumerates
force-feedback-capable
devices, and \texttt{OpenEXT(id)} / \texttt{OpenFromJoystickEXT(joystickId)} /
\texttt{OpenFromMouseEXT()} each return an owning \cnaclass{HapticDevice} handle. That handle's
own surface splits cleanly into two tiers: a simple \texttt{InitRumbleEXT()} /
\texttt{PlayRumbleEXT(strength, lengthMs)} / \texttt{StopRumbleEXT()} trio for the common
single-motor-strength case (the same conceptual level as \texttt{GamePad::SetVibration}, just
reachable for non-gamepad haptic hardware too), and a full \texttt{HapticEffectEXT}
lifecycle --- \texttt{CreateEffectEXT} / \texttt{UpdateEffectEXT} / \texttt{RunEffectEXT} /
\texttt{StopEffectEXT} / \texttt{DestroyEffectEXT}, plus \texttt{GetEffectStatusEXT} to poll
whether an effect is running. Shaped feedback such as a directional kick or sustained curve is
available only when \texttt{IsEffectSupportedEXT} accepts that effect type.

Like \cnaclass{Sensors}, this wrapper assumes its SDL subsystem already exists. No input-module
path calls \texttt{SDL\_InitSubSystem(SDL\_INIT\_HAPTIC)}, and the focused haptic tests replace
the backend rather than exercising subsystem startup. Unless the host or another module has
initialized it, enumeration/open calls can fail before hardware capability is relevant. Treat
subsystem ownership as an application integration requirement for both extension families.

\begin{lstlisting}
// A haptic-capable mouse is outside GamePad::SetVibration's scope.
if (Haptics::IsMouseHapticEXT())
{
    HapticDevice mouseHaptic = Haptics::OpenFromMouseEXT();
    if (mouseHaptic.IsOpenEXT() && mouseHaptic.InitRumbleEXT())
    {
        mouseHaptic.PlayRumbleEXT(0.5f, 150);   // half strength, 150 ms
    }
    // mouseHaptic disposes (closes the device) automatically at scope exit.
}
\end{lstlisting}

\section{Focus loss and held input}

\begin{contractnote}
Desktop focus loss and gain update \texttt{Game::IsActive} and raise the corresponding events.
Keyboard and mouse snapshots are retained across focus loss, matching FNA. Gate gameplay input on
\texttt{IsActive}; do not expect focus loss to synthesize key-up or button-up events.
\end{contractnote}

\begin{historynote}
\texttt{Game::PollEvents()} once handled only mobile background/foreground events. It now also
handles \texttt{SDL\_EVENT\_WINDOW\_FOCUS\_LOST} and \texttt{FOCUS\_GAINED}. FNA additionally
restores an X11 fullscreen-desktop flag and toggles the screensaver; CNA does not reproduce those
side effects at the pin. \texttt{WindowFocusLostDoesNotClearHeldKeysMatchingFna} preserves the
held-state policy.
\end{historynote}

\begin{lstlisting}
// Gate held input on the current focus state.
void MyGame::Update(GameTime& gameTime)
{
    if (!getIsActiveProperty())
    {
        // CNA retains keyboard and mouse snapshots while unfocused, but
        // gameplay should ignore them here.
        return;
    }

    const KeyboardState keyboard = Keyboard::GetState();
    if (keyboard.IsKeyDown(Keys::Space))
    {
        // Process held input only while the window is active.
    }
}
\end{lstlisting}

\section{Platform-specific boundaries}

A few platform constraints are external to CNA. On Wayland,
\texttt{SDL\_GetGlobalMouseState} silently returns $(0,0)$, because the compositor's own
security model forbids querying the global pointer position at all --- absolute mouse
warping is similarly focus-gated, while relative (pointer-lock) mouse mode works normally.
On Windows, XInput controllers report no USB vendor/product ID, so the \texttt{EXT}
GUID query returns \texttt{"xinput"} instead of a hexadecimal identifier.
Some platforms enumerate a touch device only after its first interaction. CNA's capability query
therefore combines live SDL enumeration with an observed-touch fallback; the separate collection
flag retains the first-touch boundary described above. In a browser via Emscripten, gamepads
remain invisible to SDL until the user
presses a button on one, a browser privacy restriction, and relative mouse mode maps onto the
Pointer Lock API, which itself requires a preceding user gesture to engage. Non-US keyboard
layouts expose an XNA limitation: several
accented and non-ASCII characters common on European layouts have no corresponding
\cnaclass{Keys} enum value at all and are simply dropped in keycode mode --- \texttt{TextInputEXT}
is the correct API for capturing that text instead.

\section{Status}

The 522-case corpus stated at the chapter opening remains bounded by fifteen formally blocked,
hardware-only campaign tasks: physical keyboard, mouse, gamepad, touchscreen, IME, and high-DPI
display checks. These are absent observations, not automated passes. The \cnaclass{Joysticks} /
\cnaclass{Sensors} / \cnaclass{Power}
\texttt{CNAEXT} extensions have fake-backend unit coverage but no demo-UI or manual-hardware
verification pass; \cnaclass{Sensors} also has the concrete subsystem-initialization gap
described above.
